\documentclass[10pt,a4paper]{amsart}
\usepackage[margin=19mm]{geometry}
\usepackage{amsmath,amssymb,mathtools}
\usepackage[scr=rsfs,cal=euler]{mathalfa}
\usepackage{xfrac}
\usepackage[unicode,hidelinks,bookmarks=false]{hyperref}
\newcommand{\ie}{i.e.\ }
\newcommand{\cf}{cf.\ }
\newcommand{\resp}{respectively\ }
\newcommand{\Subsection}{\subsection}
\providecommand{\nobreakdash}{\nobreak\hbox{-}}
\setlength{\emergencystretch}{2em}
\multlinegap0pt
\usepackage{mathtools}
\usepackage{xcolor}
\usepackage{float}
\usepackage{amssymb}
\usepackage[shortlabels]{enumitem}
\newtheorem{lemma}{Lemma}
\def\CC{\mathbb{C}}
\def\HH{\mathbb{H}}
\def\NN{\mathbb{N}}
\def\SS{\mathbb{S}}
\def\x{\chi}
\title{Orthogonal Polynomials, a Szeg\H{o}--Verblunsky Theorem and Baxter's Theorem on the Quaternionic Sphere}
\author{Connor J. Gauntlett, David P. Kimsey}
\date{Source: arXiv:2605.07787v1 (CC BY 4.0)}
\begin{document}
\maketitle

\noindent\emph{Source and licence.} Orthogonal Polynomials, a Szeg\H{o}--Verblunsky Theorem and Baxter's Theorem on the Quaternionic Sphere. Version arXiv:2605.07787v1 (CC BY 4.0), distributed by arXiv under the Creative Commons Attribution 4.0 International licence, http://creativecommons.org/licenses/by/4.0/, as recorded in the arXiv licence metadata for this exact version and shown on the abstract page https://arxiv.org/abs/2605.07787v1. Full licence notice: \texttt{licenses/arxiv-2605.07787-CC-BY-4.0.txt}.\par\smallskip

\noindent\textit{Editorial scope.} Counted unit: the article's lemma lem:UnitaryEquiv (source0.tex lines 449--455). The article's complete original proof of it is delivered with it (source0.tex lines 456--511, one \texttt{proof} environment of 56 lines). No mathematics is written, repaired or re-derived: the statement and the proof are the article's own lines, in the article's own notation, and no retained line is substituted or re-flowed.  No further source-local context is needed: every cross-reference of every retained line resolves inside the two delivered spans.  Nothing was omitted from any retained span, so the package carries no omission record.  No retained line contains a citation, so the wrapper carries no bibliography and no bibliography entry.  No retained line contains a figure environment, an image-inclusion command or drawing code, so no image is delivered and the package declares no asset dependency.  Editorial formatting only, in addition: an AMS \texttt{amsart} preamble that restates the article's own declarations and macros used by the retained lines, namely the environments \texttt{lemma} on separate counters, together with the standard packages those lines need.  The retained environments print in this document as Lemma 1 in order of appearance, whereas the article prints the same items with the numbering of its own full document.  The complete native source of the article as distributed in the arXiv e-print for this exact version is bundled unchanged as \texttt{sources/200109/source0.tex}.
\begin{lemma}
    \label{lem:UnitaryEquiv}
    Let \(i,j \in \SS\) be orthogonal. For each \(n\in\NN\), there exists some unitary \(U_n \in \CC_i^{2n \times 2n}\) with the property that, for any \(A = [a_{k\ell}]_{k,\ell=1}^n \in \HH^{n\times n}\), one has 
    \[
        \x_i(A) = U_n \begin{bmatrix} \x_i(a_{11}) & \cdots & \x_i(a_{1n}) \\ \vdots & \ddots & \vdots \\ \x_i(a_{n1}) & \cdots & \x_i(a_{nn}) \end{bmatrix} U_n^*.
    \]
\end{lemma}

\begin{proof}
    The unitary equivalence of the two matrices is in fact given by a permutation of the rows and columns, i.e. the matrix \(U_n\) is a permutation matrix. We claim that \(U_n = [u_{k\ell}]_{k,\ell=1}^{2n}\) is given by
    \[
        u_{k\ell} = \begin{cases}1, \text{ if } (k, \ell) \in \{(m, 2m-1), (n+m, 2m) : m = 1, \ldots, n\}, \\ 0, \text{ otherwise}.\end{cases}
    \]
    For example, for \(n = 2,3\) these are
    \[
        U_2 = \begin{bmatrix} 1 & 0 & 0 & 0 \\ 0 & 0 & 1 & 0 \\ 0 & 1 & 0 & 0 \\ 0 & 0 & 0 & 1 \end{bmatrix}, \quad U_3 = \begin{bmatrix} 1 & 0 & 0 & 0 & 0 & 0 \\ 0 & 0 & 1 & 0 & 0 & 0 \\ 0 & 0 & 0 & 0 & 1 & 0 \\ 0 & 1 & 0 & 0 & 0 & 0 \\ 0 & 0 & 0 & 1 & 0 & 0 \\ 0 & 0 & 0 & 0 & 0 & 1 \end{bmatrix}.
    \]
    Notice first that for any \(n \in \NN\), we may obtain \(U_{n}\) from the \(2n \times 2n\) matrix \(U_{n-1} \oplus I_2\) by moving the \((2n-1)^{\text{th}}\) row to position \(n+1\), and then shifting all following rows \(n+1, n+2, \ldots, 2n-2\) down one; in other words, we apply the permutation \(2n-1 \to n+1 \to n+2 \to \ldots \to 2n-2 \to 2n-1\). If we define the permutation matrix \(V_n\) for \(n \in \NN\) via
    \[
        V_{n+1} := I_{n} \oplus \begin{bmatrix} 0 & 1 & 0 & \cdots & 0 \\ \vdots & \ddots & \ddots & \ddots & \vdots \\ 0 & \cdots & 0 & 1 & 0 \\ 1 & 0 & \cdots & \cdots & 0 \\ 0 & \cdots & \cdots & 0 & 1 \end{bmatrix} \in \CC_i^{2(n+1) \times 2(n+1)},
    \]
    then this permutation may be described for \(n \in \NN\) by \(U_{n+1} = V_{n+1} \cdot (U_n \oplus I_2)\).
    
    Now, to see the claim, suppose for the purpose of induction that for some \(n \in \NN\) we had that
    \[
        \x_i(A) = U_n [\x_i(a_{k\ell})]_{k,\ell=1}^n U_n^*
    \]
    for all \(A = [a_{k\ell}]_{k,\ell=1}^n \in \HH^{n\times n}\). Let \(A = [a_{k\ell}]_{k,\ell=1}^{n+1} \in \HH^{(n+1)\times(n+1)}\) and calculate for an arbitrary block matrix \(\begin{bmatrix} X & Y \\ Z & W \end{bmatrix}\) that for a appropriately-sized matrices \(\Gamma\) and \(I_k\),
    \[
        \begin{bmatrix} \Gamma & 0 \\ 0 & I_k \end{bmatrix} \begin{bmatrix} X & Y \\ Z & W \end{bmatrix} \begin{bmatrix} \Gamma & 0 \\ 0 & I_k\end{bmatrix}^* = \begin{bmatrix} \Gamma X \Gamma^* & \Gamma Y \\ Z \Gamma^* & W \end{bmatrix}.
    \]
    By splitting \([\x_i(a_{k,\ell})]_{k,\ell=1}^{n+1}\) into the appropriate blocks, then, we have that
    \begin{equation}
    \label{eq:UnitaryEquivLemma}
        \begin{split}
            U_{n+1} [\x_i(a_{k,\ell})]_{k,\ell=1}^{n+1} U_{n+1}^* & = V_{n+1} \begin{bmatrix} U_n & 0 \\ 0 & I_2 \end{bmatrix} [\x_i(a_{k,\ell})]_{k,\ell=1}^{n+1} \begin{bmatrix} U_n & 0 \\ 0 & I_2 \end{bmatrix}^* V_{n+1}^* \\
            & = V_{n+1} \begin{bmatrix} U_n [\x_i(a_{k,\ell})]_{k,\ell=1}^n U_n^* & U_n [\x_i(a_{k,n+1})]_{k=1}^n \\[6pt] [\x_i(a_{n+1,\ell})]_{\ell=1}^n U_n^* & \x_i(a_{n+1}) \end{bmatrix} V_{n+1}^*. 
        \end{split}
    \end{equation}
    By the inductive hypothesis, the \((1,1)\)-block of the centre matrix is simply \(\x_i(A')\). Expand the right-most column and bottom row to write the right-hand side as
    \[
        V_{n+1} \left[\begin{array}{c|c} \x_i(A') & U_n \begin{bmatrix}a_{1,n+1}^{(1)} & a_{1,n+1}^{(2)} \\[6pt] -\overline{a_{1,n+1}^{(2)}} & \overline{a_{1,n+1}^{(1)}} \\ \vdots & \vdots \\ a_{n,n+1}^{(1)} & a_{n,n+1}^{(2)} \\[6pt] -\overline{a_{n,n+1}^{(2)}} & \overline{a_{n,n+1}^{(1)}} \end{bmatrix} \\[50pt] \hline \\[-8pt] \begin{bmatrix}a_{n+1,1}^{(1)} & a_{n+1,1}^{(2)} & \cdots & a_{n+1,n}^{(1)} & a_{n+1,n}^{(2)} \\[6pt] -\overline{a_{n+1,1}^{(2)}} & \overline{a_{n+1,1}^{(1)}} & \cdots & -\overline{a_{n+1,n}^{(2)}} & \overline{a_{n+1,n}^{(1)}}\end{bmatrix} U_n^* & \begin{array}{cc} a_{n+1,n+1}^{(1)} & a_{n+1,n+1}^{(2)} \\[6pt] -\overline{a_{n+1,n+1}^{(2)}} & \overline{a_{n+1, n+1}^{(1)}}\end{array} \end{array}\right] V_{n+1}^*.
    \]
    Of course, multiplication by \(U_n\) on the left and by \(U_n^*\) on the right simply permutes the rows and columns, respectively, so applying these permutations we obtain the matrix
    \[
        V_{n+1} \left[\begin{array}{c|c} \x_i(A') & \begin{array}{cc} a_{1,n+1}^{(1)} & a_{1,n+1}^{(2)} \\[6pt] \vdots & \vdots \\ a_{n,n+1}^{(1)} & a_{n,n+1}^{(2)} \\[6pt] -\overline{a_{1,n+1}^{(2)}} & \overline{a_{1,n+1}^{(1)}} \\ \vdots & \vdots \\ -\overline{a_{n,n+1}^{(2)}} & \overline{a_{n,n+1}^{(1)}} \end{array} \\[10pt] \\[-8pt] \hline \\[-8pt] \begin{array}{cccccc} a_{n+1,1}^{(1)} & \cdots &  a_{n+1,n}^{(1)} & a_{n+1,1}^{(2)} & \cdots & a_{n+1,n}^{(2)} \\[6pt] -\overline{a_{n+1,1}^{(2)}} & \cdots & -\overline{a_{n+1,n}^{(2)}} & \overline{a_{n+1,1}^{(1)}} & \cdots & \overline{a_{n+1,n}^{(1)}}\end{array} & \begin{array}{cc} a_{n+1,n+1}^{(1)} & a_{n+1,n+1}^{(2)} \\[6pt] -\overline{a_{n+1,n+1}^{(2)}} & \overline{a_{n+1, n+1}^{(1)}}\end{array} \end{array}\right] V_{n+1}^*.
    \]

    Now, we have described above the permutation on rows arising from multiplying by the matrix \(V_{n+1}\) on the left; multiplying by \(V_{n+1}^*\) on the right performs the same permutation on the columns. Thus --- applying first the permutation on the rows --- our matrix becomes
    \[
        \begin{bmatrix} a_{11}^{(1)} & \cdots & a_{1n}^{(1)} & a_{11}^{(2)} & \cdots & a_{1n}^{(2)} & a_{1,n+1}^{(1)} & a_{1,n+1}^{(2)} \\[6pt] \vdots & \ddots & \vdots & \vdots & \ddots & \vdots & \vdots & \vdots \\[6pt] a_{n1}^{(1)} & \cdots & a_{nn}^{(1)} & a_{n1}^{(2)} & \cdots & a_{nn}^{(2)}  & a_{n,n+1}^{(1)} & a_{n,n+1}^{(2)} \\[6pt] a_{n+1,1}^{(1)} & \cdots & a_{n+1,n}^{(1)} & a_{n+1,1}^{(2)} & \cdots & a_{n+1,n}^{(2)}& a_{n+1,n+1}^{(1)} & a_{n+1,n+1}^{(2)} \\[6pt] -\overline{a_{11}^{(2)}} & \cdots & -\overline{a_{1n}^{(2)}} & \overline{a_{11}^{(1)}} & \cdots & \overline{a_{1n}^{(1)}} & -\overline{a_{1,n+1}^{(2)}} & \overline{a_{1,n+1}^{(1)}} \\[6pt] \vdots & \ddots & \vdots & \vdots & \ddots & \vdots & \vdots & \vdots \\[6pt] -\overline{a_{n1}^{(2)}} & \cdots & -\overline{a_{nn}^{(2)}} & \overline{a_{n1}^{(1)}} & \cdots & \overline{a_{nn}^{(1)}} & -\overline{a_{n,n+1}^{(2)}} & \overline{a_{n,n+1}^{(1)}} \\[6pt] -\overline{a_{n+1,1}^{(2)}} & \cdots & -\overline{a_{n+1,n}^{(2)}} & \overline{a_{n+1,1}^{(1)}} & \cdots & \overline{a_{n+1,n}^{(1)}} & -\overline{a_{n+1,n+1}^{(2)}} & \overline{a_{n+1,n+1}^{(1)}} \end{bmatrix} V_{n+1}^*,
    \]
    and finally, by permuting the columns in the same manner we arrive at
    \[
        \begin{bmatrix} a_{11}^{(1)} & \cdots & a_{1n}^{(1)} & a_{1,n+1}^{(1)} &  a_{11}^{(2)} & \cdots & a_{1n}^{(2)} & a_{1,n+1}^{(2)} \\[6pt] \vdots & \ddots & \vdots & \vdots & \vdots & \ddots & \vdots & \vdots \\[6pt] a_{n1}^{(1)} & \cdots & a_{nn}^{(1)} & a_{n,n+1}^{(1)} & a_{n1}^{(2)} & \cdots & a_{nn}^{(2)} & a_{n,n+1}^{(2)} \\[6pt] a_{n+1,1}^{(1)} & \cdots & a_{n+1,n}^{(1)} & a_{n+1,n+1}^{(1)} & a_{n+1,1}^{(2)} & \cdots & a_{n+1,n}^{(2)} & a_{n+1,n+1}^{(2)} \\[6pt] -\overline{a_{11}^{(2)}} & \cdots & -\overline{a_{1n}^{(2)}} & -\overline{a_{1,n+1}^{(2)}} & \overline{a_{11}^{(1)}} & \cdots & \overline{a_{1n}^{(1)}} & \overline{a_{1,n+1}^{(1)}} \\[6pt] \vdots & \ddots & \vdots & \vdots & \vdots & \ddots & \vdots & \vdots \\[6pt] -\overline{a_{n1}^{(2)}} & \cdots & -\overline{a_{nn}^{(2)}} & -\overline{a_{n,n+1}^{(2)}} & \overline{a_{n1}^{(1)}} & \cdots & \overline{a_{nn}^{(1)}} & \overline{a_{n,n+1}^{(1)}} \\[6pt] -\overline{a_{n+1,1}^{(2)}} & \cdots & -\overline{a_{n+1,n}^{(2)}} & -\overline{a_{n+1,n+1}^{(2)}} & \overline{a_{n+1,1}^{(1)}} & \cdots & \overline{a_{n+1,n}^{(1)}} & \overline{a_{n+1,n+1}^{(1)}} \end{bmatrix} = \x_i(A).
    \]
    Hence the right-hand side of \eqref{eq:UnitaryEquivLemma} may now be readily seen to be \(\x_i(A)\).

    It remains to note the base case: when \(n=1\), \(A = a_{11} \in \HH\) and furthermore \(U_1 = I_2\), so
    \[
        U_1 \x_i(a_{11}) U_1^* = \x_i(a_{11}) = \x_i(A).
    \]
    By induction, then, \(U_n [\x_i(a_{k\ell})]_{k,\ell=1}^n U_n^* = \x_i(A)\) for all \(n \in \NN\) and \(A = [a_{k\ell}]_{k,\ell=1}^n \in \HH^{n\times n}\).
\end{proof}

\end{document}
